\chapter{C1 UH: a homogeneous biquadratic conductor}
\label{component:pf7ffa921793c}
\noindent\textbf{Component }\path{pf7ffa921793c}\par
{\small\noindent Source: \path{R13/advances/c1_uh/uh_all_tests.tex}\par}
{\small\noindent SHA256: \path{f7ffa921793c0b1b432463bba3a4de95659a1186ce3299ca64731620b8b2ed33}\par}
{\small\noindent Inventory use: UH1--UH20 and complete source-boundary finite-cap proof | C1\_\allowbreak{}BARE\_\allowbreak{}OTHER\_\allowbreak{}TERMINALS; precise subordinate uses retained below\par}
\medskip
\section{C1 UH: a homogeneous biquadratic conductor}
\label{pf7ffa921793c:sec:r7-c1-uh}
All coefficients below are divided by the original positive pure-source
$d$. We treat precisely (NI9) and (NI10), with no alteration of the native
source (NC1)--(NC7). The revision-6 common-event theorem already covers
all tests except positive unit $u$ and $w=1/v$ infinitesimal.
Put $x=k^2$, $K=1+x$, $\sigma=1/\eta$, $z=1/\zeta=q^2$,
$b=r/z$, $H=h/\sigma$, and retain $\sigma>0$ (it need not be a unit).
On (NI9), $\sigma$ is infinitesimal and $z$ is a unit; on (NI10),
$z$ is infinitesimal and $\sigma$ is bounded. Throughout, $k,H$ are
positive units, $b,z,\sigma$ are bounded, $Q=1-xb>0$, and
\[
 r=zb,\quad R=1+xzb,\quad M=R-z+\sigma zH,
 \quad J=1+\sigma(xH-1+zb),\quad\theta=1-\sigma H,
 \quad E_s=1-xH.
\]
The exact source polynomials, all with bounded arguments, are
\begin{align*}
 P={}&xK+\sigma x+\sigma z+\sigma zb x^2
 +x^2\sigma zH(1-\sigma H)+2x\sigma^2zH
 +\sigma^2z^2bx^2H,\\
 V={}&z(1+x\sigma H)E_s+x(1-zb)(1+\sigma zH),\\
 B_\alpha={}&R\sigma zH+RxH+ME_s.
\end{align*}
Writing $m=m_{\rm old}/z$, $c=C/z$, $a=Az$, $\ell=\alpha z$,
\[
 c={K(1+\sigma zb)E_s\over2(1-zb)P},\quad
 m={Kx(1+\sigma zb)H(\sigma z+x)\over2MP},\quad
 a={V\over2bP},\quad
 \ell={x^2(1+\sigma zb)B_\alpha\over2QP}.                 \tag{UH1}
\]
Define the polynomial quotients
\[
 N_D={K(\sigma z+x)(1+\sigma zb)-JP\over\sigma},\qquad
 N_\beta={Kx(1+\sigma zb)(1+\sigma zH)-\theta P\over\sigma}.
\]
Then $D=\sigma N_D/(2JP)$ and
$\beta=\sigma N_\beta/(2\theta P)$, exactly.

\paragraph{The bounded full parent.}
Let $D_u=1+u^2$, $D_w=1+w^2$. Its six event amplitudes are
\[
 M_w=W-wX+uY-uwZ,\ R_w=uW-Y,\ Y,\
 J_w=uwW+uX-wY-Z,\ Y-wZ,\ B_w=X+uZ.
\]
Thus the exact parent is
\[
 F=4mzM_w^2+4czD_wR_w^2+4(a/z)D_uD_wY^2
 +4DJ_w^2+4(\ell/z)D_u(Y-wZ)^2+4\beta D_wB_w^2+2K_w, \tag{UH2}
\]
where, without any restriction to a state ray,
\begin{align*}
 K_w={}&u^2w^2Z^2-2u^2wYZ-2uwWZ+2uw^2XZ
       +u^2w^2W^2+2u^2wWX-2uw^2WY\\
      &+(D_u+w^2)(X^2+Y^2)-2uwXY.
\end{align*}
For a potentially unbounded weight choose a finite cap, keeping bounded
weights unchanged. Subtract and retain the following nonnegative
centered excesses:
\begin{align*}
 &4(m-m_*)[qM_w-kM(wZ/q)]^2,\\
 &4(c-c_*)D_w[qR_w-x(1-r)(wZ/q)]^2,\\
 &4(a-a_*)D_uD_w[(Y-xbwZ)/q]^2,
 \quad4(\ell-\ell_*)D_u[(Y-xbwZ)/q]^2.          \tag{UH3}
\end{align*}
The first line uses $qM_w$ and the second $qR_w$; thus these are
exactly excesses of the scaled weights in (UH1). Their remainders
are bounded because $mM,c(1-r),ab,\ell Q$ are bounded.
When $\sigma$ is a unit, additionally retain
\[
 4(D-D_*)(J_w+JZ)^2,
 \qquad4(\beta-\beta_*)D_w(X+uhZ)^2.               \tag{UH4}
\]
Here $DJ$ and $\beta\theta$ are bounded. When $\sigma$ is
infinitesimal, no (UH4) cap is needed: $D,\beta=O(\sigma^2)$,
as proved below. Denote the complete remainder by $G$.
Every retained row vanishes at the actual contact
\[
 u=k,\quad w={q\sqrt\sigma\over k},\qquad
 (W,X,Y,Z)=({\sqrt\sigma\over q},-kh,kbq\sqrt\sigma,1).
                                                               \tag{UH5}
\]
No comparison between independent source and test rates occurs.

\paragraph{Exact homogeneous conductor, including every first normal.}
Put $h=\sigma H$, $\delta=u-k$, and use the six roots
\[
\begin{gathered}
 N_x=X+khZ,\quad N_y=(Y-kbq^2W)/q,\\
 U=wW-\sigma Z/k,\quad N=(wZ-q^2W/k)/q,\\
 S=q\delta W,\quad T=\sigma\delta Z.
\end{gathered}                                     \tag{UH6}
\]
Substitute $X=N_x-khZ$, $Y=qN_y+kbq^2W$ in the entire
bounded remainder $G$. Its complete quadratic normal block is the
coefficient collection in $(N_x,N_y)$. Its scalar and first-normal
coefficients have exactly the following supports:
\begin{align*}
 G_{00}&={A\over q^2}w^2Z^2+B w^2W^2+C\sigma^2 Z^2
                  +D_1 wWZ+E q^2W^2,\\
 G_{10}&=A_xwW+B_xZ+C_xw^2Z,\\
 G_{01}&=q(A_y+B_yw^2)W+C_ywZ/q.                    \tag{UH7}
\end{align*}
All eleven displayed coefficient families and the quadratic normal
coefficients are bounded polynomials in $u$, with bounded source
coefficients; $B_x$ is divisible by $\sigma$. In the unit-$\sigma$
case divisibility means a bounded quotient, not a new polynomial
identity. Formulas below specify every coefficient before caps, and
(UH3)--(UH4) specify their exact changes. No state normal is omitted.

Set $s=\sigma/k$, $a_0=q^2/k$ and
\[
 L=2A/k+2B\sigma/k+D_1,\quad
 J_0=C-B/k^2,\quad K_0=A/k^2+2B\sigma/k^2+D_1/k+E.
\]
The scalar is identically
\[
 A N^2+B U^2+qLWN+\sigma^2J_0Z^2+q^2K_0W^2.       \tag{UH8}
\]
Exact source contact and all stationary jets give
\[
 L(k)=J_0(k)=K_0(k)=0,\qquad K_0'(k)=-\sigma J_0'(k),
\]
\[
 B_x(k)+s[A_x(k)+a_0C_x(k)]=0,\qquad
 A_y(k)+C_y(k)/k+s a_0B_y(k)=0.                    \tag{UH9}
\]
These hold separately for every centered excess, hence for every cap.
Write the monic divisions
\[
 L=\delta L_1,\quad J_0=\delta j+\delta^2J_2,\quad
 K_0=-\sigma\delta j+\delta^2K_2,
\]
\[
 B_x+s(A_x+a_0C_x)=\sigma\delta H_x,\qquad
 A_y+C_y/k+s a_0B_y=\delta H_y.
\]
All quotients are bounded. The scalar and full first normals become
\begin{align*}
 G_{00}={}&A N^2+B U^2+(L_1+k\sigma j)SN-kjTU
                          +J_2T^2+K_2S^2,\\
 G_{10}={}&(A_x+a_0C_x)U+C_xqwN+H_xT,\\
 G_{01}={}&(C_y+s q^2B_y)N+qB_ywU+H_yS.            \tag{UH10}
\end{align*}
Together with the unchanged full quadratic normal block, this is an
exact bounded six-root Gram for $G$, not a scalar restriction.
The only root relation needed for strict completion is
\[
 T^2=\sigma S^2+k\sigma\delta SN-k\delta TU.        \tag{UH11}
\]
Its verification is one multiplication using (UH6). Adding any bounded
multiple of its zero changes no parent. No mixed inverse, root, order,
or source/test identification occurs in (UH6)--(UH11).

For the cap argument it is useful to print every centered row directly
in this same frame. In the order $m,c,(a,\ell),D,\beta$, they are
\begin{align*}
 &uzN_y-(kR+z\delta)N-qwN_x-(zQ/k)S,\\
 &S-zN_y-x(1-r)N,\qquad N_y-xbN,\\
 &(u-kbz)U+uN_x-qwN_y+(E_s/k)T,\qquad N_x+HT.
                                                               \tag{UH20}
\end{align*}
These are exact identities before taking any residue.

\paragraph{Explicit coefficient input.}
For (UH2) before caps put $j_u=u-kbz$, $p_u=1+ukbz$,
$t_u=kh-u$. The full scalar coefficients in (UH7) are
\begin{align*}
 A={}&4mz^2t_u^2+4\ell D_u+(4\beta+2)zt_u^2,\\
 B={}&(4cz+4D+2)j_u^2+4aD_uxb^2z,\\
 C={}&[4D(1+ukh)^2+4\beta(u-kh)^2+2D_uxh^2]/\sigma^2,\\
 D_1={}&8mzp_ut_u-8Dj_u(1+ukh)-8\ell D_ukb
          -4u^2k(h+bz)-4u+4uxhbz,\\
 E={}&4mp_u^2+4cj_u^2+4(a+\ell)D_uxb^2+2D_uxb^2z.
\end{align*}
The first normals are
\begin{align*}
 A_x={}&-8mzp_u+(8D+4)uj_u,\\
 B_x={}&-8Du(1+ukh)+8\beta(u-kh)-4D_ukh,\\
 C_x={}&-8mzt_u+(8\beta+4)(u-kh),\\
 A_y={}&8mzup_u-8czj_u+8(a+\ell)D_ukb+4D_ukbz,\\
 B_y={}&-8czj_u+8aD_ukb-(8D+4)j_u,\\
 C_y={}&8mz^2ut_u-8\ell D_u+8Dz(1+ukh)+4uzt_u.
\end{align*}
The coefficients of $N_x^2,N_y^2,N_xN_y$ are respectively
\[
 4mzw^2+4Du^2+4\beta D_w+2(D_u+w^2),
\]
\[
 4mz^2u^2+4cz^2D_w+4aD_uD_w+4\ell D_u+4Dzw^2
                            +2z(D_u+w^2),
 \qquad-8mzuqw-8Duqw-4uqw.                         \tag{UH12}
\]
Here products such as $Du$ mean $D\,u$, whereas $D_u=1+u^2$.

\paragraph{Infinitesimal $\sigma$: source collapse and strict residue.}
The exact polynomial $P_m=(N_D+xN_\beta)/\sigma$ is bounded.
At $\sigma=0$, $N_D=-xN_\beta$ and
\[
 N_\beta=x(K+z)H+bxz-x-z.
\]
Nonnegativity of $D,\beta$ forces
\[
 H=H_0+O(\sigma),\quad
 H_0={x+zQ\over x(K+z)},\qquad E_0=1-xH_0={R\over K+z}.
                                                               \tag{UH13}
\]
Indeed $JD+x\theta\beta=\sigma^2P_m/(2P)$ and $J,\theta$
are units. Thus $D,\beta=O(\sigma^2)$, including either literal zero.
Moreover
\[
 (D+x\beta)/\sigma^2\longrightarrow KH_0E_0/2>0.
\]
The other limiting coefficients are
\[
 c_0={R\over2x(1-zb)(K+z)},\quad m_0={xH_0\over2(1-zQ)},
 \quad a_0^{\rm wt}={H_0\over2b},\quad
 \ell_0={xR\over2Q(K+z)}.                           \tag{UH14}
\]
At a matched angle $u=k$, the residue Gram (UH10) splits into
$(N_y,N,S)$ and $(N_x,U,T)$, because their only off-block coefficients
contain $w$. Before caps, the $(N_y,N)$ block is positive definite;
its angular Schur complement is
\[
 \Gamma=8RH_0/K>0.                                \tag{UH15}
\]
The $(N_x,U)$ block is precisely
\[
 2K(N_x+xH_0U)^2+2KxH_0E_0U^2.                   \tag{UH16}
\]
The remaining $T$ couplings are bounded. As $\sigma,\delta$ have
zero residue, (UH11) permits an arbitrarily large positive $T^2$
pivot without changing any other residue coefficient. This proves a
strict full six-root residue Gram. For $u^2\ne x$, use the full raw
frame $(E_t,F_t,G_t,R_t,X,T_0)=(qW,Y/q,wZ/q,wW,X,\sigma Z)$.
Its leading first block is
\begin{align*}
 &4m_0[E_t+uz(F_t-G_t)]^2+4c_0(uE_t-zF_t)^2
       +4a_0^{\rm wt}D_uF_t^2+4\ell_0D_u(F_t-G_t)^2\\
 &\qquad+2z\{u^2(G_t-F_t)^2+F_t^2\}-4uE_tG_t.
\end{align*}
Put $v_0=u^2$ and $M_0=1-zQ$. In $(F_t,F_t-G_t)$ the hard block
is diagonal, with weights
\[
 B_h={2(x+zQ)[v_0(1-r)+R]\over bx(1-r)(K+z)},\qquad
 L_h={2R[v_0(x+zQ)+xM_0]\over Q(K+z)M_0}.
\]
The exact $E_t$ Schur complement is
\[
 {2(v_0-x)^2(1+v_0)RH_0\over
 [v_0(1-r)+R][v_0(x+zQ)+xM_0]}.                  \tag{UH17}
\]
It has a positive-unit limit through every opening when $v_0\ne x$.
The hard minimizers are
\[
 {F_t\over E_t}={ubK\over v_0(1-r)+R},\qquad
 {F_t-G_t\over E_t}={-uKQ\over v_0(x+zQ)+xM_0},
\]
also bounded through every opening. The other raw blocks are
$2\{u^2R_t^2+2u^2R_tX+D_uX^2\}$ and
$4(D+u^2\beta)T_0^2/\sigma^2$, both strictly positive.
This gives the full bounded away-angle lift, without a determinant-only
claim or any omitted first state normal.

\paragraph{Unit $\sigma$: angular Schur sign.}
Here (NI10) gives $z=0$ at residue. Put $h=\sigma H$,
$J=1+\sigma(xH-1)$, $\theta=1-\sigma H$. The source signs imply
\[
 0<\sigma<K,\qquad {1\over K+\sigma}\le H
       \le{x+\sigma\over x(K+\sigma)}.             \tag{UH18}
\]
These inequalities are stated on the interior; its closure includes
$\sigma=K$ when $J=\theta=0$.
The limits are
\[
 m={KxH\over2(K+\sigma)},\quad c={K(1-xH)\over2x(K+\sigma)},
 \quad a={1\over2b(K+\sigma)},\quad
 \ell={x\over2Q(K+\sigma)},
\]
\[
 D={K\over2J(K+\sigma)}-\tfrac12,
 \quad\beta={K\over2\theta(K+\sigma)}-\tfrac12.
\]
The two hard blocks in (UH10) are positive definite. Their angular
Schur coefficients $\gamma_S,\gamma_T$ have no $ST$ term and satisfy
\[
 \gamma_S+\sigma\gamma_T
 ={ -2p(H)\over\sigma(K+\sigma)}>0,\qquad
 p(H)=Kx(K+4\sigma)H^2-2K(x+2\sigma)H+x.           \tag{UH19}
\]
Convexity of $p$ and (UH18) prove strict negativity: its endpoint
values are respectively
\[
 {-\sigma(3\sigma x+4\sigma+4x+4)\over(K+\sigma)^2},\quad
 {-\sigma(4\sigma x+3\sigma+4x^2+4x)\over x(K+\sigma)^2}.
\]
Use (UH11) to distribute this positive total between $S^2,T^2$;
for example choose its multiple so both resulting Schur coefficients
are positive. This yields the required strict six-root Gram.

For completeness, angles $u^2\ne x$ are strictly positive as well.
Eliminating the two hard state variables from the leading form,
multiplying by the positive source factor $(K+\sigma)/2$, and
writing $U_0=u^2$, $P_0=(w/q)(Z/q)$, $T_0=(Z/q)^2$ gives a positive
Laurent function
\[
 A+B P_0^2+C P_0^2/T_0+D_0T_0-2E P_0,
\]
where
\[
 A=KxH+K(1-xH)U_0/x,\ B=x(1+U_0),\
 C={KU_0\over J+\theta U_0},
\]
\[
 D_0=(1+U_0)\left[{K(1+U_0)\over J+\theta U_0}-(K+\sigma)\right],
 \quad E={K\sqrt{U_0}(1+U_0)\over J+\theta U_0}.
\]
The strict criterion is $E<\sqrt{AB}+\sqrt{CD_0}$.
Squaring only to analyze these positive real residue quantities gives
\[
 (E^2-AB-CD_0)^2-4ABCD_0
 ={K^2(1+U_0)^2\over(J+\theta U_0)^2}(U_0-x)^2Q_2(U_0).
\]
The quadratic $Q_2$ has positive constant $H^2x^2J^2$, positive
leading coefficient $\theta^2(1-xH)^2$, and $Q_2(x)=xp(H)<0$.
Explicitly, with $E_s=1-xH$,
\[
 Q_2(U_0)=E_s^2\theta^2U_0^2
 +2E_s\{x\theta^2+E_s\theta J+(K+\sigma)\theta-2K\}U_0
 +x^2H^2J^2.
\]
The concave quadratic
$U_0(K+\sigma)-[x^2H+(1-xH)U_0](J+\theta U_0)$
is negative at zero and infinity and positive at $x$; at each of its
two zeros the preceding discriminant is negative. Consequently its
positive interval lies inside $Q_2<0$. This proves the criterion for
all $U_0\ne x$. Sign charts with $P_0\le0$ are immediate. These
square roots are an analysis of real residues only and are not tensor
roots or mixed localizations.

This inequality is directly a bounded Gram allocation, not a
comparison of mixed ratios. Before choosing projective coordinates,
put $E_t=qW$, $G_t=wZ/q$, $R_t=wW$. The exact relation
$E_tG_t=R_tZ$ permits distributing the cross term between the two
quadratic blocks $(E_t,G_t)$ and $(R_t,Z)$. The strict inequality
above is exactly the condition that both $2\times2$ blocks can be
made positive definite by a finite real allocation. It therefore
handles every independent rate of $q,w$.
At $J=0$ or $\theta=0$, the root-interval argument permits one endpoint
at zero or infinity. At the double face $J=\theta=0$, necessarily
$\sigma=K$, $H=1/K$. As $J+\theta U_0\to0$,
\[
 E-\sqrt{CD_0}\longrightarrow K\sqrt{U_0},\qquad
 AB-K^2U_0=(U_0-x)^2.
\]
Thus the constrained allocation remains strictly positive away from
$U_0=x$, including this double face without dividing by a vanishing
residue denominator.

\paragraph{Opening faces, tensor conclusion, and exact scope.}
The finite-cap argument is applied to the full Gram above. For
$\sigma$ infinitesimal the possible diverging rows have openings
$1-zb$, $1-zQ$, $b$, $Q$; at most two can diverge together. The pairs
$(c,m)$, $(m,a)$, $(c,\ell)$ have independent hard rows; other pairs
are excluded by $0<zb,xb<1$ and $1-zQ>0$. In the $(N_y,N)$ columns
of (UH20), their determinants have absolute values $kzK,k,x$,
respectively; each is a unit on its compatible face. The hard block
is uniformly coercive, since it contains positive-unit multiples
of $N_y^2$ and $(N_y-N)^2$. Its minimizing hard coordinates for
fixed angular $S$ are explicitly
\[
 {N_y\over S}={b(2bxz-x+1)\over K},\qquad
 {N\over S}={4bxz-x-2z+1\over xK}.
\]
They are bounded through every opening. Thus the limit of the
uncapped Schur complement is the Schur complement restricted to the
diverging centered rows' kernel. The quotients (UH15)--(UH16) stay
positive units there. The elementary finite-dimensional penalty
lemma now applies: a bounded symmetric form strictly positive on
the common kernel becomes strictly positive after sufficiently large
finite positive multiples of these rows' squares are added. This
proves the finite-cap claim, with the entire angular right-hand sides
in (UH20) retained. Away from $u^2=x$, (UH17) and its printed bounded
hard minimizers give the same constrained strictness directly.

For unit $\sigma$, only $b,Q,J,\theta$ can be openings.
The $b,Q$ rows coincide in the hard direction $N_y-xbN$ and cannot
both diverge. The $J,\theta$ rows on $(N_x,U)$ are independent even
if both vanish; their determinant is the unit $u$. Formula (UH19)
remains a positive unit through these faces. More explicitly, the
two hard blocks dominate positive-unit multiples of
$N_y^2+(N_y-N)^2$ and $N_x^2+(N_x+U)^2$. Their angular minimizers are
\[
 {N_y\over S}={b(1-HxK)\over K},\quad
 {N\over S}={1-HxK\over xK},
\]
\[
 {N_x\over T}={\sigma H^2xK-HxK-\sigma H(2x+1)+x\over\sigma K},
 \quad {U\over T}={HxK(K+\sigma)-xK-\sigma\over x\sigma K}.
\]
Their $D$ and $\beta$ row defects divided by $T$ are respectively
$k(HK-1)J/(\sigma K)$ and $-x\theta(HK-1)/(\sigma K)$.
Hence every relevant constrained limit exists with bounded
minimizers; the same penalty argument applies. On the angular
relation $T^2=\sigma S^2$, its strictly positive Schur value is
(UH19), and after taking finite caps the bounded zero (UH11)
distributes that value into two strict angular pivots. The away-angle
allocation and its single/double-face limits above give the same
conclusion there. All removed excesses remain exact positive
source/test squares. No positive lower bound on an original opening
has been imposed.

A strict bounded residue Gram lifts by the inherited tensor Gram
lemma. Every root in (UH6) is a finite sum of separated source and
testing factors. After restoring the retained squares this certifies
the complete original parent, including $D=0$ or $\beta=0$.
Together with revision-6 remote tests, this closes both (NI9) and
(NI10). Those are the entire actual UH portion of C1.5. The already
proved UU and UI portions and exact source stabilizers therefore close
the entire C1.5 row, subject to the exact checks and mathematical review
recorded alongside this section.
